链表和指针重连深度题卡 按题型拆成章内大节，但每个大节都先从 LeetCode 案例切入，再进入分流、案例矩阵、案例推导和实验复盘。这样读者看到的是从具体题推到规律，而不是先读抽象方法再猜它能用在哪。

\topic{链表与指针重连}
这一节先看 LeetCode 2 两数相加、LeetCode 19 删除链表的倒数第 N 个结点、LeetCode 21 合并两个有序链表 这几道 LeetCode 题，再整理同一题型的分流和推导。阅读顺序不是先背原理，而是先看案例的暴力瓶颈，再把重复出现的观察抽象成方法。

\subsection{原理和适用条件}

以 LeetCode 2 两数相加、LeetCode 19 删除链表的倒数第 N 个结点、LeetCode 21 合并两个有序链表 为主线：LeetCode 2 两数相加：模型是“链表按低位到高位保存数字，本质是竖式加法而不是整数转换。”，优化抓手是“维护两个游标和进位，任一链未结束或进位非零都继续生成节点。”；LeetCode 19 删除链表的倒数第 N 个结点：模型是“倒数位置可以转成两个指针之间固定 n 步的距离。”，优化抓手是“快指针先走 n 步，再让快慢指针同步移动，慢指针停在待删节点前驱。”；LeetCode 21 合并两个有序链表：模型是“两个链表已经各自有序，下一节点只能来自两个当前头中较小者。”，优化抓手是“虚拟头统一结果链表，尾指针不断接上较小节点并推进来源链。”。这些案例共同推出的规律是：先用数组或多次扫描得到直观答案，再把操作改成节点重连。链表题的关键是保存后继、明确前驱和当前段边界，不能靠值交换掩盖指针关系。 方法名只是复盘后的标签，真正要掌握的是从题面约束、暴力失败、状态设计到边界反例的推导链。

\begin{principlebox}{图示：从 LeetCode 案例到 链表与指针重连推导链}
\begin{tabular}{@{}p{0.18\linewidth}p{0.74\linewidth}@{}}
暴力基线 & 枚举候选、完整模拟或逐个检查，先保证覆盖全部可能答案。\\
瓶颈定位 & 链表慢点在于不能随机访问。要找倒数位置、分组边界或交点，只能通过指针间距离、快慢指针或哈希记录访问过的节点。\\
结构选择 & 优化来自哨兵节点、快慢指针、前驱指针和局部反转。链表题的核心不是值交换，而是保持断开和接回的顺序正确。\\
不变量 & 必须明确每个指针指向哪一段：已经处理好的前缀头尾、当前待处理节点、下一段头节点。每次改 \code{next} 前先保存后继。\\
代码落地 & 节点通常由评测平台管理，代码不要随意 \code{delete} 输入节点。若创建新节点，要说明所有权。比较交点时比较地址，不比较值。\\
\end{tabular}
\end{principlebox}

\subsection{LeetCode 案例分流}

\begin{principlebox}{从 LeetCode 案例分流：链表与指针重连}
\textbf{基础重连。}
代表案例：LeetCode 206 反转链表、LeetCode 21 合并两个有序链表、LeetCode 19 删除倒数第 N 个节点。先看这些题的题面条件：题目要求反转、合并、删除或重排节点。由此推出的处理方式是：使用虚拟头、前驱、当前和后继指针维护链结构。风险点：每次改 next 前保存后继，防止断链。

\textbf{快慢指针。}
代表案例：LeetCode 141 环形链表、LeetCode 142 环形链表 II、LeetCode 234 回文链表。先看这些题的题面条件：题目涉及中点、倒数位置、环入口或回文结构。由此推出的处理方式是：用两个速度或间距不同的指针制造位置关系。风险点：比较交点和环入口时比较节点地址，不比较值。

\textbf{分组和局部反转。}
代表案例：LeetCode 24 两两交换链表节点、LeetCode 25 K 个一组翻转链表、LeetCode 92 反转链表 II。先看这些题的题面条件：题目按 K 个、两两或某个区间反转链表。由此推出的处理方式是：先找到组边界，保存下一组头，再做局部反转和接回。风险点：不足一组是否反转要按题意处理。

\textbf{多链表和稳定节点。}
代表案例：LeetCode 23 合并 K 个升序链表、LeetCode 148 排序链表、LeetCode 160 相交链表。先看这些题的题面条件：多个链表需要合并、排序或交叉定位。由此推出的处理方式是：用堆、分治、双指针或自底向上归并维护节点身份。风险点：不要用值相等代替节点相同。

\end{principlebox}

\subsection{案例矩阵}

本节案例覆盖 LeetCode 2 两数相加、LeetCode 19 删除链表的倒数第 N 个结点、LeetCode 21 合并两个有序链表、LeetCode 24 两两交换链表中的节点、LeetCode 25 K 个一组翻转链表、LeetCode 61 旋转链表、LeetCode 82 删除排序链表中的重复元素 II、LeetCode 83 删除排序链表中的重复元素、LeetCode 92 反转链表 II、LeetCode 141 环形链表、LeetCode 142 环形链表 II、LeetCode 146 LRU 缓存、LeetCode 148 排序链表、LeetCode 160 相交链表、LeetCode 206 反转链表、LeetCode 234 回文链表、LeetCode 430 扁平化多级双向链表、LeetCode 460 LFU 缓存。阅读时不要按题号背答案，而要先判断题面落在哪个分流场景：查询目标是什么，输入约束给了什么单调性或等价关系，暴力慢在哪里，优化结构保存了什么状态。

\begin{principlebox}{案例矩阵}
\textbf{LeetCode 2 两数相加。}
链表按低位到高位保存数字，本质是竖式加法而不是整数转换。单点风险：两链长度不同、最后进位、输入为零、不能用内置整数承载超长数字。

\textbf{LeetCode 19 删除链表的倒数第 N 个结点。}
倒数位置可以转成两个指针之间固定 n 步的距离。单点风险：删除头节点、n 等于链表长度、只有一个节点、题目保证 n 合法。

\textbf{LeetCode 21 合并两个有序链表。}
两个链表已经各自有序，下一节点只能来自两个当前头中较小者。单点风险：某一条链为空、重复值、剩余整段接回、是否复用原节点。

\textbf{LeetCode 24 两两交换链表中的节点。}
链表按固定长度为二的小组做局部重连。单点风险：空链、单节点、奇数长度、交换后前驱更新到组尾。

\textbf{LeetCode 25 K 个一组翻转链表。}
链表被切成长度为 K 的连续组，只有完整组才反转。单点风险：不足 K 个不反转、K 为一、刚好整除、反转后头尾接错。

\textbf{LeetCode 61 旋转链表。}
右旋 k 位等价于把链表尾部一段切到头部。单点风险：空链、单节点、k 为零、k 大于长度、刚好整圈。

\textbf{LeetCode 82 删除排序链表中的重复元素 II。}
有序链表中重复值连续出现，题目要求把所有重复值整组删除。单点风险：头部重复、尾部重复、全部重复、没有重复。

\textbf{LeetCode 83 删除排序链表中的重复元素。}
有序链表让重复节点相邻，只需保留每段相同值的第一个节点。单点风险：空链、单节点、全重复、重复段在尾部。

\textbf{LeetCode 92 反转链表 II。}
只反转链表中一段连续区间，区间外节点必须保持原连接。单点风险：从头开始反转、只反转一个节点、反转到尾部、区间边界。

\textbf{LeetCode 141 环形链表。}
快慢指针在有环时必然相遇，无环时快指针先到空。单点风险：空链、单节点自环、两节点环、尾部无环。

\textbf{LeetCode 142 环形链表 II。}
相遇点到入口的距离关系来自快慢指针速度差。单点风险：入口在头节点、长尾短环、无环、单节点环。

\textbf{LeetCode 146 LRU 缓存。}
哈希表负责按 key 定位，双向链表负责维护新旧顺序。单点风险：容量为一、重复 put、get 不存在、更新已有值。

\textbf{LeetCode 148 排序链表。}
链表不能随机访问，适合归并排序而不是快速排序数组 partition。单点风险：空链、单节点、重复值、奇数长度、切断子链。

\textbf{LeetCode 160 相交链表。}
相交是节点地址相同，不是节点值相同。单点风险：不相交、头节点相交、尾节点相交、长度差很大。

\textbf{LeetCode 206 反转链表。}
每次把当前节点的 next 指向已经反转好的前缀头。单点风险：空链、单节点、两个节点、不要丢失剩余链。

\textbf{LeetCode 234 回文链表。}
链表回文要求前半段和后半段逆序后逐一相等。单点风险：空链、单节点、奇数长度、偶数长度、是否恢复链表。

\textbf{LeetCode 430 扁平化多级双向链表。}
多级链表按深度优先顺序展开成一条双向链表。单点风险：child 为空、next 和 child 同时存在、多层嵌套、prev 指针恢复。

\textbf{LeetCode 460 LFU 缓存。}
缓存淘汰同时按访问频次和同频最近性排序，需要维护两个维度的索引。单点风险：容量为零、更新已有 key、频次链表变空时最小频次更新、同频淘汰最旧节点。

\end{principlebox}

\subsection{共性落点}

\begin{principlebox}{本组共性落点}
\textbf{慢版本。} 暴力做法常把链表拷到数组里处理，再写回或重建。这样虽然降低指针难度，却失去原地节点重连的训练价值。
\textbf{瓶颈。} 链表慢点在于不能随机访问。要找倒数位置、分组边界或交点，只能通过指针间距离、快慢指针或哈希记录访问过的节点。
\textbf{状态字段。} 虚拟头、上一段尾、当前段头、当前节点、后继节点；任何改 next 的操作之前都要保存后续入口。
\textbf{不变量。} 必须明确每个指针指向哪一段：已经处理好的前缀头尾、当前待处理节点、下一段头节点。每次改 \code{next} 前先保存后继。
\textbf{实现检查。} 所有重连步骤按保存后继、断开或反转、接回前缀、接回后缀的顺序写；最后检查是否形成环或丢节点。
\textbf{C++ 落地。} 节点通常由评测平台管理，代码不要随意 \code{delete} 输入节点。若创建新节点，要说明所有权。比较交点时比较地址，不比较值。
\textbf{证明抓手。} 证明从链结构开始：已处理链连通、未处理链可达、二者连接点明确，节点没有丢失也没有成环。
\end{principlebox}

\subsection{案例推导}

\noindent\textbf{LeetCode 2 两数相加。}

\textbf{题目说明。} LeetCode 2 两数相加 的题面入口要先还原成具体任务：链表按低位到高位保存数字，本质是竖式加法而不是整数转换。

\textbf{样例入口。} 拿 2->4->3 加 5->6->4 手算，个位 2+5 得到 7，十位 4+6 得到 10，所以当前位写 0 并把进位 1 带到下一位。再走百位 3+4+1 得到 8。

\textbf{暴力基线。} 慢版本可以先把节点顺序抄到数组里，按数组推导目标顺序。回到链表后，难点不再是值的顺序，而是节点身份和 next 指针不能丢。放到本题就是：先按“链表按低位到高位保存数字，本质是竖式加法而不是整数转换”完整试慢版本；样例里已经能看到“规律是链表从低位开始，状态只需要两个游标、当前进位和结果尾节点；两条链长度不同或最后还有进位时，循环仍不能停”，所以优化前必须先能说清慢版本枚举了哪些候选。

\textbf{暴力过程。} 拿 2->4->3 加 5->6->4 手算，个位 2+5 得到 7，十位 4+6 得到 10，所以当前位写 0 并把进位 1 带到下一位。再走百位 3+4+1 得到 8。

\textbf{重复工作。} 题面模型：链表按低位到高位保存数字，本质是竖式加法而不是整数转换。暴力版的慢点要落到本题：慢版本会围绕“链表按低位到高位保存数字，本质是竖式加法而不是整数转换”借助数组或多次扫描确认目标顺序。样例规律是“规律是链表从低位开始，状态只需要两个游标、当前进位和结果尾节点；两条链长度不同或最后还有进位时，循环仍不能停”，说明优化点不在节点值，而在少量指针角色能否保持已处理链、未处理链和接回位置都不丢。后面的优化只消掉这类重复工作，不能改变答案集合。

\textbf{优化条件。} 本题的关键观察是：维护两个游标和进位，任一链未结束或进位非零都继续生成节点。这句话必须能翻译成可维护状态；如果题目条件改变后观察不再成立，就不能继续套原来的写法。

\textbf{相邻状态。} 规律是链表从低位开始，状态只需要两个游标、当前进位和结果尾节点；两条链长度不同或最后还有进位时，循环仍不能停。把这个特例继续拆开看：每个指针都要有角色，上一段尾、当前段头、当前节点和后继入口不能混。只要一次改 next 之前没有保存后继，就可能丢掉整段未处理链。

\textbf{状态复用。} 把数组辅助结构换成几个稳定指针角色；重连顺序必须保证已处理链合法、未处理链仍可达。本题落点是：规律是链表从低位开始，状态只需要两个游标、当前进位和结果尾节点；两条链长度不同或最后还有进位时，循环仍不能停。手算时只改被新输入影响的那一小部分，而不是回头重算完整候选。

\textbf{指针和字段。} 状态字段要能说清：虚拟头、上一段尾、当前段头、当前节点、后继节点；任何改 next 的操作之前都要保存后续入口。手算时按这个协议跟踪：拿 2->4->3 加 5->6->4 手算，个位 2+5 得到 7，十位 4+6 得到 10，所以当前位写 0 并把进位 1 带到下一位。再走百位 3+4+1 得到 8。然后按协议复查：手算时画出 prev、cur、next 和下一段头节点。每次改 next 前先保存后继，这是链表推导的安全线。

\textbf{代码落点。} 所有重连步骤按保存后继、断开或反转、接回前缀、接回后缀的顺序写；最后检查是否形成环或丢节点。关键代码行必须能逐句翻译成“旧状态怎样变成新状态”，否则就只是模板。

\textbf{正确性。} 证明从链结构开始：已处理链连通、未处理链可达、二者连接点明确，节点没有丢失也没有成环。

\textbf{边界实验。} 先用两链长度不同、最后进位、输入为零、不能用内置整数承载超长数字做反例检查。实验时覆盖两链长度不同、最后进位、输入为零、不能用内置整数承载超长数字，每步画出前驱、当前、后继和下一段头。链表题不要只看输出值，还要检查节点是否丢失或成环。

\textbf{复杂度变化。} 暴力可能多次扫描链表或拷贝到数组；优化后每个节点通常只被访问、重连常数次，目标是线性时间和常数额外空间。

\textbf{迁移追问。} 如果题目改成“若链表高位在前，可以用栈反向处理或先反转链表”，先判断原状态是否仍然覆盖新答案集合，再决定是微调边界、改 value 语义，还是换算法。

\noindent\textbf{LeetCode 19 删除链表的倒数第 N 个结点。}

\textbf{题目说明。} LeetCode 19 删除链表的倒数第 N 个结点 的题面入口要先还原成具体任务：倒数位置可以转成两个指针之间固定 n 步的距离。

\textbf{样例入口。} 拿长度 5 删除倒数第 2 个手算。快指针先走 2 步，慢指针和快指针保持固定距离；当快指针到尾后，慢指针正好停在待删节点前驱。再试删除头节点，若没有虚拟头就缺前驱。

\textbf{暴力基线。} 暴力先扫描一遍链表求长度 len，再计算要删除的是第 len-n+1 个节点；第二遍走到它的前驱后改 next。

\textbf{暴力过程。} 以 1->2->3->4->5、n=2 为例，暴力先数出长度 5，目标是正数第 4 个节点 4，再从头走到节点 3，把 3->next 改成 5。

\textbf{重复工作。} 题面模型：倒数位置可以转成两个指针之间固定 n 步的距离。暴力版的慢点要落到本题：浪费在为了定位倒数位置先完整数长度。倒数第 n 个节点和它前驱，可以通过快慢指针之间固定 n+1 步距离一次扫描定位。后面的优化只消掉这类重复工作，不能改变答案集合。

\textbf{优化条件。} 本题的关键观察是：快指针先走 n 步，再让快慢指针同步移动，慢指针停在待删节点前驱。这句话必须能翻译成可维护状态；如果题目条件改变后观察不再成立，就不能继续套原来的写法。

\textbf{相邻状态。} 规律是虚拟头统一头部删除，快慢指针固定距离定位前驱。把这个特例继续拆开看：每个指针都要有角色，上一段尾、当前段头、当前节点和后继入口不能混。只要一次改 next 之前没有保存后继，就可能丢掉整段未处理链。

\textbf{状态复用。} 使用虚拟头 dummy 指向 head。fast 先从 dummy 走 n+1 步，然后 fast 和 slow 同步走；fast 到空时，slow 正好是待删节点前驱。手算时只改被新输入影响的那一小部分，而不是回头重算完整候选。

\textbf{指针和字段。} 状态字段要能说清：dummy 负责统一删除头节点，fast 负责制造距离，slow 停在前驱，slow->next 是要删除的节点。手算时按这个协议跟踪：删除 1->2->3->4->5 的倒数第 2 个时，fast 先领先 slow 三步；同步移动到 fast 为空，slow 在 3，删除 slow->next=4。

\textbf{代码落点。} 关键代码是 dummy.next=head，fast=\&dummy，预走 n+1 步；最后 slow->next=slow->next->next。预走从 dummy 开始，才能覆盖删除头节点。关键代码行必须能逐句翻译成“旧状态怎样变成新状态”，否则就只是模板。

\textbf{正确性。} fast 和 slow 始终相隔 n+1 条边；当 fast 走到链表末尾之后，slow 后面正好还有 n 个节点，所以 slow 是目标前驱。

\textbf{边界实验。} 先用删除头节点、n 等于链表长度、只有一个节点、题目保证 n 合法做反例检查。实验时覆盖删除头节点、n 等于链表长度、只有一个节点、题目保证 n 合法，每步画出前驱、当前、后继和下一段头。链表题不要只看输出值，还要检查节点是否丢失或成环。

\textbf{复杂度变化。} 两遍扫描 O(n)；快慢指针仍是 O(n) 时间，但只走一遍主体扫描，额外空间 O(1)。

\textbf{迁移追问。} 如果题目改成“若 n 不保证合法，需要在快指针预走阶段检测长度不足”，先判断原状态是否仍然覆盖新答案集合，再决定是微调边界、改 value 语义，还是换算法。

\noindent\textbf{LeetCode 21 合并两个有序链表。}

\textbf{题目说明。} LeetCode 21 合并两个有序链表 的题面入口要先还原成具体任务：两个链表已经各自有序，下一节点只能来自两个当前头中较小者。

\textbf{样例入口。} 拿 1->3 和 2->4 手算，先接 1，再比较 3 和 2 接 2，之后接 3 和 4。

\textbf{暴力基线。} 暴力可以把两个链表的值全部放入数组，排序后重新建链表。

\textbf{暴力过程。} 以 1->3 和 2->4 为例，数组法得到 [1, 2, 3, 4]；但原链表本来有序，每一步只需要比较两个当前头。

\textbf{重复工作。} 题面模型：两个链表已经各自有序，下一节点只能来自两个当前头中较小者。暴力版的慢点要落到本题：浪费在完整排序和新建节点。两个链表各自有序，下一节点一定来自 l1 或 l2 的较小头节点。后面的优化只消掉这类重复工作，不能改变答案集合。

\textbf{优化条件。} 本题的关键观察是：虚拟头统一结果链表，尾指针不断接上较小节点并推进来源链。这句话必须能翻译成可维护状态；如果题目条件改变后观察不再成立，就不能继续套原来的写法。

\textbf{相邻状态。} 规律是结果尾指针每次接较小头节点，某条链耗尽后直接接另一条剩余部分；虚拟头统一处理首节点。把这个特例继续拆开看：每个指针都要有角色，上一段尾、当前段头、当前节点和后继入口不能混。只要一次改 next 之前没有保存后继，就可能丢掉整段未处理链。

\textbf{状态复用。} 用 dummy 作为结果虚拟头，tail 指向结果尾。每次比较 l1 和 l2，接上较小节点并推进对应链表。手算时只改被新输入影响的那一小部分，而不是回头重算完整候选。

\textbf{指针和字段。} 状态字段要能说清：l1/l2 是两条未处理链的当前头，tail 是已合并链表尾部，dummy.next 是最终结果入口。手算时按这个协议跟踪：1 和 2 比较接 1；3 和 2 比较接 2；3 和 4 比较接 3；最后 l1 空，直接接剩余 4。

\textbf{代码落点。} 关键是 tail->next=chosen 后 tail=tail->next，再推进 chosen 来源链。某条链为空后 tail->next 接另一条剩余链。关键代码行必须能逐句翻译成“旧状态怎样变成新状态”，否则就只是模板。

\textbf{正确性。} 两条链都已排序，两个当前头中较小者不可能被另一条链后面的节点超过，所以它就是全局下一个节点。

\textbf{边界实验。} 先用某一条链为空、重复值、剩余整段接回、是否复用原节点做反例检查。实验时覆盖某一条链为空、重复值、剩余整段接回、是否复用原节点，每步画出前驱、当前、后继和下一段头。链表题不要只看输出值，还要检查节点是否丢失或成环。

\textbf{复杂度变化。} 排序数组法 O((m+n)log(m+n)) 且额外空间高；双指针合并 O(m+n) 时间，O(1) 额外空间。

\textbf{迁移追问。} 如果题目改成“合并 K 条链表可用分治或堆，核心仍是维护每条链的当前头”，先判断原状态是否仍然覆盖新答案集合，再决定是微调边界、改 value 语义，还是换算法。

\noindent\textbf{LeetCode 24 两两交换链表中的节点。}

\textbf{题目说明。} LeetCode 24 两两交换链表中的节点 的题面入口要先还原成具体任务：链表按固定长度为二的小组做局部重连。

\textbf{样例入口。} 拿 1->2->3->4 手算，第一组交换后应是 2->1，并且 1 要接到下一组交换后的头。

\textbf{暴力基线。} 慢版本可以先把节点顺序抄到数组里，按数组推导目标顺序。回到链表后，难点不再是值的顺序，而是节点身份和 next 指针不能丢。放到本题就是：先按“链表按固定长度为二的小组做局部重连”完整试慢版本；样例里已经能看到“规律是每次固定 pre、first、second、next 四个指针，先保存 next，再重连 second->first”，所以优化前必须先能说清慢版本枚举了哪些候选。

\textbf{暴力过程。} 拿 1->2->3->4 手算，第一组交换后应是 2->1，并且 1 要接到下一组交换后的头。

\textbf{重复工作。} 题面模型：链表按固定长度为二的小组做局部重连。暴力版的慢点要落到本题：慢版本会围绕“链表按固定长度为二的小组做局部重连”借助数组或多次扫描确认目标顺序。样例规律是“规律是每次固定 pre、first、second、next 四个指针，先保存 next，再重连 second->first”，说明优化点不在节点值，而在少量指针角色能否保持已处理链、未处理链和接回位置都不丢。后面的优化只消掉这类重复工作，不能改变答案集合。

\textbf{优化条件。} 本题的关键观察是：使用组前驱、第一节点、第二节点和下一组头，按顺序改指针并推进前驱。这句话必须能翻译成可维护状态；如果题目条件改变后观察不再成立，就不能继续套原来的写法。

\textbf{相邻状态。} 规律是每次固定 pre、first、second、next 四个指针，先保存 next，再重连 second->first，first->next。把这个特例继续拆开看：每个指针都要有角色，上一段尾、当前段头、当前节点和后继入口不能混。只要一次改 next 之前没有保存后继，就可能丢掉整段未处理链。

\textbf{状态复用。} 把数组辅助结构换成几个稳定指针角色；重连顺序必须保证已处理链合法、未处理链仍可达。本题落点是：规律是每次固定 pre、first、second、next 四个指针，先保存 next，再重连 second->first。手算时只改被新输入影响的那一小部分，而不是回头重算完整候选。

\textbf{指针和字段。} 状态字段要能说清：虚拟头、上一段尾、当前段头、当前节点、后继节点；任何改 next 的操作之前都要保存后续入口。手算时按这个协议跟踪：拿 1->2->3->4 手算，第一组交换后应是 2->1，并且 1 要接到下一组交换后的头。然后按协议复查：手算时画出 prev、cur、next 和下一段头节点。每次改 next 前先保存后继，这是链表推导的安全线。

\textbf{代码落点。} 所有重连步骤按保存后继、断开或反转、接回前缀、接回后缀的顺序写；最后检查是否形成环或丢节点。关键代码行必须能逐句翻译成“旧状态怎样变成新状态”，否则就只是模板。

\textbf{正确性。} 证明从链结构开始：已处理链连通、未处理链可达、二者连接点明确，节点没有丢失也没有成环。

\textbf{边界实验。} 先用空链、单节点、奇数长度、交换后前驱更新到组尾做反例检查。实验时覆盖空链、单节点、奇数长度、交换后前驱更新到组尾，每步画出前驱、当前、后继和下一段头。链表题不要只看输出值，还要检查节点是否丢失或成环。

\textbf{复杂度变化。} 暴力可能多次扫描链表或拷贝到数组；优化后每个节点通常只被访问、重连常数次，目标是线性时间和常数额外空间。

\textbf{迁移追问。} 如果题目改成“K 个一组翻转是同一思想的推广，只是组内反转更长”，先判断原状态是否仍然覆盖新答案集合，再决定是微调边界、改 value 语义，还是换算法。

\noindent\textbf{LeetCode 25 K 个一组翻转链表。}

\textbf{题目说明。} LeetCode 25 K 个一组翻转链表 的题面入口要先还原成具体任务：链表被切成长度为 K 的连续组，只有完整组才反转。

\textbf{样例入口。} 拿 1->2->3->4、k=2 画图。第一组边界是 1 和 2，下一组入口是 3；反转前如果不保存 3，链会断。反转后上一组尾接 2，旧头 1 接 3。

\textbf{暴力基线。} 暴力可以把每 K 个节点的值放进数组反转，再写回链表；这避开了指针重连，但没有真正按节点反转。

\textbf{暴力过程。} 以 1->2->3->4->5、k=2 为例，第一组 1、2 反转成 2->1，第二组 3、4 反转成 4->3，剩余 5 不动。

\textbf{重复工作。} 题面模型：链表被切成长度为 K 的连续组，只有完整组才反转。暴力版的慢点要落到本题：浪费在数组辅助和节点值交换。题目要求链表节点重排时，应直接反转每一组的 next 指针。后面的优化只消掉这类重复工作，不能改变答案集合。

\textbf{优化条件。} 本题的关键观察是：每轮先探测 K 个节点确认完整，再保存下一组头并做局部反转接回。这句话必须能翻译成可维护状态；如果题目条件改变后观察不再成立，就不能继续套原来的写法。

\textbf{相邻状态。} 规律是每组先找 kth 和 group\_next，再局部反转并接回前后两段。把这个特例继续拆开看：每个指针都要有角色，上一段尾、当前段头、当前节点和后继入口不能混。只要一次改 next 之前没有保存后继，就可能丢掉整段未处理链。

\textbf{状态复用。} dummy 指向头。group\_prev 是上一组尾，先从 group\_prev 出发找第 k 个节点 group\_end；不足 k 个就停止。保存 next\_group 后反转 [group\_start, group\_end]，再接回。手算时只改被新输入影响的那一小部分，而不是回头重算完整候选。

\textbf{指针和字段。} 状态字段要能说清：group\_prev 是已处理部分尾，group\_start 是本组头，group\_end 是本组尾，next\_group 是下一组入口。手算时按这个协议跟踪：第一组 start=1、end=2、next\_group=3。反转后得到 2->1，再让 group\_prev->next=2，1->next=3。

\textbf{代码落点。} 关键顺序是先保存 next\_group，再局部反转；反转结束后原 group\_start 变成本组尾，下一轮 group\_prev=group\_start。关键代码行必须能逐句翻译成“旧状态怎样变成新状态”，否则就只是模板。

\textbf{正确性。} 每次只改变完整 K 组内部指针，并把本组头尾接回已处理链和未处理链；不足 K 个节点不进入反转，符合题意。

\textbf{边界实验。} 先用不足 K 个不反转、K 为一、刚好整除、反转后头尾接错做反例检查。实验时覆盖不足 K 个不反转、K 为一、刚好整除、反转后头尾接错，每步画出前驱、当前、后继和下一段头。链表题不要只看输出值，还要检查节点是否丢失或成环。

\textbf{复杂度变化。} 数组辅助 O(n) 空间；原地分组反转每节点访问常数次，O(n) 时间、O(1) 空间。

\textbf{迁移追问。} 如果题目改成“递归写法短但可能有栈风险，迭代写法更接近工程实现”，先判断原状态是否仍然覆盖新答案集合，再决定是微调边界、改 value 语义，还是换算法。

\noindent\textbf{LeetCode 61 旋转链表。}

\textbf{题目说明。} LeetCode 61 旋转链表 的题面入口要先还原成具体任务：右旋 k 位等价于把链表尾部一段切到头部。

\textbf{样例入口。} 拿 1->2->3->4->5、k=2 手算，右移后应为 4->5->1->2->3；先连成环，再从新尾 3 处断开。

\textbf{暴力基线。} 慢版本可以先把节点顺序抄到数组里，按数组推导目标顺序。回到链表后，难点不再是值的顺序，而是节点身份和 next 指针不能丢。放到本题就是：先按“右旋 k 位等价于把链表尾部一段切到头部”完整试慢版本；样例里已经能看到“规律是 k 先对长度取模，新头是第 n-k 个节点的后继，断环前必须保存新头”，所以优化前必须先能说清慢版本枚举了哪些候选。

\textbf{暴力过程。} 拿 1->2->3->4->5、k=2 手算，右移后应为 4->5->1->2->3；先连成环，再从新尾 3 处断开。

\textbf{重复工作。} 题面模型：右旋 k 位等价于把链表尾部一段切到头部。暴力版的慢点要落到本题：慢版本会围绕“右旋 k 位等价于把链表尾部一段切到头部”借助数组或多次扫描确认目标顺序。样例规律是“规律是 k 先对长度取模，新头是第 n-k 个节点的后继，断环前必须保存新头”，说明优化点不在节点值，而在少量指针角色能否保持已处理链、未处理链和接回位置都不丢。后面的优化只消掉这类重复工作，不能改变答案集合。

\textbf{优化条件。} 本题的关键观察是：先求长度并让 k 对长度取模，再找到新尾节点并重连成新头。这句话必须能翻译成可维护状态；如果题目条件改变后观察不再成立，就不能继续套原来的写法。

\textbf{相邻状态。} 规律是 k 先对长度取模，新头是第 n-k 个节点的后继，断环前必须保存新头。把这个特例继续拆开看：每个指针都要有角色，上一段尾、当前段头、当前节点和后继入口不能混。只要一次改 next 之前没有保存后继，就可能丢掉整段未处理链。

\textbf{状态复用。} 把数组辅助结构换成几个稳定指针角色；重连顺序必须保证已处理链合法、未处理链仍可达。本题落点是：规律是 k 先对长度取模，新头是第 n-k 个节点的后继，断环前必须保存新头。手算时只改被新输入影响的那一小部分，而不是回头重算完整候选。

\textbf{指针和字段。} 状态字段要能说清：虚拟头、上一段尾、当前段头、当前节点、后继节点；任何改 next 的操作之前都要保存后续入口。手算时按这个协议跟踪：拿 1->2->3->4->5、k=2 手算，右移后应为 4->5->1->2->3；先连成环，再从新尾 3 处断开。然后按协议复查：手算时画出 prev、cur、next 和下一段头节点。每次改 next 前先保存后继，这是链表推导的安全线。

\textbf{代码落点。} 所有重连步骤按保存后继、断开或反转、接回前缀、接回后缀的顺序写；最后检查是否形成环或丢节点。关键代码行必须能逐句翻译成“旧状态怎样变成新状态”，否则就只是模板。

\textbf{正确性。} 证明从链结构开始：已处理链连通、未处理链可达、二者连接点明确，节点没有丢失也没有成环。

\textbf{边界实验。} 先用空链、单节点、k 为零、k 大于长度、刚好整圈做反例检查。实验时覆盖空链、单节点、k 为零、k 大于长度、刚好整圈，每步画出前驱、当前、后继和下一段头。链表题不要只看输出值，还要检查节点是否丢失或成环。

\textbf{复杂度变化。} 暴力可能多次扫描链表或拷贝到数组；优化后每个节点通常只被访问、重连常数次，目标是线性时间和常数额外空间。

\textbf{迁移追问。} 如果题目改成“若本来把链表接成环，最后必须断开新尾的 next”，先判断原状态是否仍然覆盖新答案集合，再决定是微调边界、改 value 语义，还是换算法。

\noindent\textbf{LeetCode 82 删除排序链表中的重复元素 II。}

\textbf{题目说明。} LeetCode 82 删除排序链表中的重复元素 II 的题面入口要先还原成具体任务：有序链表中重复值连续出现，题目要求把所有重复值整组删除。

\textbf{样例入口。} 拿 1->2->3->3->4 手算，两个 3 都要删除，不是保留一个；pre 停在 2，cur 扫完整段重复 3 后接到 4。

\textbf{暴力基线。} 慢版本可以先把节点顺序抄到数组里，按数组推导目标顺序。回到链表后，难点不再是值的顺序，而是节点身份和 next 指针不能丢。放到本题就是：先按“有序链表中重复值连续出现，题目要求把所有重复值整组删除”完整试慢版本；样例里已经能看到“规律是虚拟头加 pre 指向已确认无重复前缀尾，遇到重复段时整段跳过”，所以优化前必须先能说清慢版本枚举了哪些候选。

\textbf{暴力过程。} 拿 1->2->3->3->4 手算，两个 3 都要删除，不是保留一个；pre 停在 2，cur 扫完整段重复 3 后接到 4。

\textbf{重复工作。} 题面模型：有序链表中重复值连续出现，题目要求把所有重复值整组删除。暴力版的慢点要落到本题：慢版本会围绕“有序链表中重复值连续出现，题目要求把所有重复值整组删除”借助数组或多次扫描确认目标顺序。样例规律是“规律是虚拟头加 pre 指向已确认无重复前缀尾，遇到重复段时整段跳过”，说明优化点不在节点值，而在少量指针角色能否保持已处理链、未处理链和接回位置都不丢。后面的优化只消掉这类重复工作，不能改变答案集合。

\textbf{优化条件。} 本题的关键观察是：虚拟头保存结果前驱，遇到一段重复值时跳过整段，否则前驱前进。这句话必须能翻译成可维护状态；如果题目条件改变后观察不再成立，就不能继续套原来的写法。

\textbf{相邻状态。} 规律是虚拟头加 pre 指向已确认无重复前缀尾，遇到重复段时整段跳过。把这个特例继续拆开看：每个指针都要有角色，上一段尾、当前段头、当前节点和后继入口不能混。只要一次改 next 之前没有保存后继，就可能丢掉整段未处理链。

\textbf{状态复用。} 把数组辅助结构换成几个稳定指针角色；重连顺序必须保证已处理链合法、未处理链仍可达。本题落点是：规律是虚拟头加 pre 指向已确认无重复前缀尾，遇到重复段时整段跳过。手算时只改被新输入影响的那一小部分，而不是回头重算完整候选。

\textbf{指针和字段。} 状态字段要能说清：虚拟头、上一段尾、当前段头、当前节点、后继节点；任何改 next 的操作之前都要保存后续入口。手算时按这个协议跟踪：拿 1->2->3->3->4 手算，两个 3 都要删除，不是保留一个；pre 停在 2，cur 扫完整段重复 3 后接到 4。然后按协议复查：手算时画出 prev、cur、next 和下一段头节点。每次改 next 前先保存后继，这是链表推导的安全线。

\textbf{代码落点。} 所有重连步骤按保存后继、断开或反转、接回前缀、接回后缀的顺序写；最后检查是否形成环或丢节点。关键代码行必须能逐句翻译成“旧状态怎样变成新状态”，否则就只是模板。

\textbf{正确性。} 证明从链结构开始：已处理链连通、未处理链可达、二者连接点明确，节点没有丢失也没有成环。

\textbf{边界实验。} 先用头部重复、尾部重复、全部重复、没有重复做反例检查。实验时覆盖头部重复、尾部重复、全部重复、没有重复，每步画出前驱、当前、后继和下一段头。链表题不要只看输出值，还要检查节点是否丢失或成环。

\textbf{复杂度变化。} 暴力可能多次扫描链表或拷贝到数组；优化后每个节点通常只被访问、重连常数次，目标是线性时间和常数额外空间。

\textbf{迁移追问。} 如果题目改成“若只保留一个重复值，就是删除重复元素的较简单版本”，先判断原状态是否仍然覆盖新答案集合，再决定是微调边界、改 value 语义，还是换算法。

\noindent\textbf{LeetCode 83 删除排序链表中的重复元素。}

\textbf{题目说明。} LeetCode 83 删除排序链表中的重复元素 的题面入口要先还原成具体任务：有序链表让重复节点相邻，只需保留每段相同值的第一个节点。

\textbf{样例入口。} 拿 1->1->2 手算，当前节点 1 的 next 值相同，就让当前节点跳过 next；遇到 2 时再前进。

\textbf{暴力基线。} 慢版本可以先把节点顺序抄到数组里，按数组推导目标顺序。回到链表后，难点不再是值的顺序，而是节点身份和 next 指针不能丢。放到本题就是：先按“有序链表让重复节点相邻，只需保留每段相同值的第一个节点”完整试慢版本；样例里已经能看到“规律是有序链表让重复值相邻，只需比较当前节点和下一个节点，保留每组第一个”，所以优化前必须先能说清慢版本枚举了哪些候选。

\textbf{暴力过程。} 拿 1->1->2 手算，当前节点 1 的 next 值相同，就让当前节点跳过 next；遇到 2 时再前进。

\textbf{重复工作。} 题面模型：有序链表让重复节点相邻，只需保留每段相同值的第一个节点。暴力版的慢点要落到本题：慢版本会围绕“有序链表让重复节点相邻，只需保留每段相同值的第一个节点”借助数组或多次扫描确认目标顺序。样例规律是“规律是有序链表让重复值相邻，只需比较当前节点和下一个节点，保留每组第一个”，说明优化点不在节点值，而在少量指针角色能否保持已处理链、未处理链和接回位置都不丢。后面的优化只消掉这类重复工作，不能改变答案集合。

\textbf{优化条件。} 本题的关键观察是：当前节点和后继值相同就跳过后继，否则当前节点前进。这句话必须能翻译成可维护状态；如果题目条件改变后观察不再成立，就不能继续套原来的写法。

\textbf{相邻状态。} 规律是有序链表让重复值相邻，只需比较当前节点和下一个节点，保留每组第一个。把这个特例继续拆开看：每个指针都要有角色，上一段尾、当前段头、当前节点和后继入口不能混。只要一次改 next 之前没有保存后继，就可能丢掉整段未处理链。

\textbf{状态复用。} 把数组辅助结构换成几个稳定指针角色；重连顺序必须保证已处理链合法、未处理链仍可达。本题落点是：规律是有序链表让重复值相邻，只需比较当前节点和下一个节点，保留每组第一个。手算时只改被新输入影响的那一小部分，而不是回头重算完整候选。

\textbf{指针和字段。} 状态字段要能说清：虚拟头、上一段尾、当前段头、当前节点、后继节点；任何改 next 的操作之前都要保存后续入口。手算时按这个协议跟踪：拿 1->1->2 手算，当前节点 1 的 next 值相同，就让当前节点跳过 next；遇到 2 时再前进。然后按协议复查：手算时画出 prev、cur、next 和下一段头节点。每次改 next 前先保存后继，这是链表推导的安全线。

\textbf{代码落点。} 所有重连步骤按保存后继、断开或反转、接回前缀、接回后缀的顺序写；最后检查是否形成环或丢节点。关键代码行必须能逐句翻译成“旧状态怎样变成新状态”，否则就只是模板。

\textbf{正确性。} 证明从链结构开始：已处理链连通、未处理链可达、二者连接点明确，节点没有丢失也没有成环。

\textbf{边界实验。} 先用空链、单节点、全重复、重复段在尾部做反例检查。实验时覆盖空链、单节点、全重复、重复段在尾部，每步画出前驱、当前、后继和下一段头。链表题不要只看输出值，还要检查节点是否丢失或成环。

\textbf{复杂度变化。} 暴力可能多次扫描链表或拷贝到数组；优化后每个节点通常只被访问、重连常数次，目标是线性时间和常数额外空间。

\textbf{迁移追问。} 如果题目改成“和删除重复元素 II 的区别是是否整段删除”，先判断原状态是否仍然覆盖新答案集合，再决定是微调边界、改 value 语义，还是换算法。

\noindent\textbf{LeetCode 92 反转链表 II。}

\textbf{题目说明。} LeetCode 92 反转链表 II 的题面入口要先还原成具体任务：只反转链表中一段连续区间，区间外节点必须保持原连接。

\textbf{样例入口。} 拿 1->2->3->4->5、left=2、right=4 手算，反转中段后是 1->4->3->2->5；反转前要保存 left 前驱和 right 后继。

\textbf{暴力基线。} 慢版本可以先把节点顺序抄到数组里，按数组推导目标顺序。回到链表后，难点不再是值的顺序，而是节点身份和 next 指针不能丢。放到本题就是：先按“只反转链表中一段连续区间，区间外节点必须保持原连接”完整试慢版本；样例里已经能看到“规律是虚拟头找到子段前驱，局部头插或三指针反转，最后接回前后两段”，所以优化前必须先能说清慢版本枚举了哪些候选。

\textbf{暴力过程。} 拿 1->2->3->4->5、left=2、right=4 手算，反转中段后是 1->4->3->2->5；反转前要保存 left 前驱和 right 后继。

\textbf{重复工作。} 题面模型：只反转链表中一段连续区间，区间外节点必须保持原连接。暴力版的慢点要落到本题：慢版本会围绕“只反转链表中一段连续区间，区间外节点必须保持原连接”借助数组或多次扫描确认目标顺序。样例规律是“规律是虚拟头找到子段前驱，局部头插或三指针反转，最后接回前后两段”，说明优化点不在节点值，而在少量指针角色能否保持已处理链、未处理链和接回位置都不丢。后面的优化只消掉这类重复工作，不能改变答案集合。

\textbf{优化条件。} 本题的关键观察是：用虚拟头找到区间前驱，再用头插法把区间内部节点逐个移到前面。这句话必须能翻译成可维护状态；如果题目条件改变后观察不再成立，就不能继续套原来的写法。

\textbf{相邻状态。} 规律是虚拟头找到子段前驱，局部头插或三指针反转，最后接回前后两段。把这个特例继续拆开看：每个指针都要有角色，上一段尾、当前段头、当前节点和后继入口不能混。只要一次改 next 之前没有保存后继，就可能丢掉整段未处理链。

\textbf{状态复用。} 把数组辅助结构换成几个稳定指针角色；重连顺序必须保证已处理链合法、未处理链仍可达。本题落点是：规律是虚拟头找到子段前驱，局部头插或三指针反转，最后接回前后两段。手算时只改被新输入影响的那一小部分，而不是回头重算完整候选。

\textbf{指针和字段。} 状态字段要能说清：虚拟头、上一段尾、当前段头、当前节点、后继节点；任何改 next 的操作之前都要保存后续入口。手算时按这个协议跟踪：拿 1->2->3->4->5、left=2、right=4 手算，反转中段后是 1->4->3->2->5；反转前要保存 left 前驱和 right 后继。然后按协议复查：手算时画出 prev、cur、next 和下一段头节点。每次改 next 前先保存后继，这是链表推导的安全线。

\textbf{代码落点。} 所有重连步骤按保存后继、断开或反转、接回前缀、接回后缀的顺序写；最后检查是否形成环或丢节点。关键代码行必须能逐句翻译成“旧状态怎样变成新状态”，否则就只是模板。

\textbf{正确性。} 证明从链结构开始：已处理链连通、未处理链可达、二者连接点明确，节点没有丢失也没有成环。

\textbf{边界实验。} 先用从头开始反转、只反转一个节点、反转到尾部、区间边界做反例检查。实验时覆盖从头开始反转、只反转一个节点、反转到尾部、区间边界，每步画出前驱、当前、后继和下一段头。链表题不要只看输出值，还要检查节点是否丢失或成环。

\textbf{复杂度变化。} 暴力可能多次扫描链表或拷贝到数组；优化后每个节点通常只被访问、重连常数次，目标是线性时间和常数额外空间。

\textbf{迁移追问。} 如果题目改成“K 个一组反转是在多个固定长度区间上重复这一类局部重连”，先判断原状态是否仍然覆盖新答案集合，再决定是微调边界、改 value 语义，还是换算法。

\noindent\textbf{LeetCode 141 环形链表。}

\textbf{题目说明。} LeetCode 141 环形链表 的题面入口要先还原成具体任务：快慢指针在有环时必然相遇，无环时快指针先到空。

\textbf{样例入口。} 拿 1->2->3 且 3 指回 2 手算，快指针每次走两步，慢指针每次走一步，进入环后二者距离会被不断缩小直到相遇。

\textbf{暴力基线。} 慢版本可以先把节点顺序抄到数组里，按数组推导目标顺序。回到链表后，难点不再是值的顺序，而是节点身份和 next 指针不能丢。放到本题就是：先按“快慢指针在有环时必然相遇，无环时快指针先到空”完整试慢版本；样例里已经能看到“规律是快慢指针的相遇说明有环；无环时快指针会先到空”，所以优化前必须先能说清慢版本枚举了哪些候选。

\textbf{暴力过程。} 拿 1->2->3 且 3 指回 2 手算，快指针每次走两步，慢指针每次走一步，进入环后二者距离会被不断缩小直到相遇。

\textbf{重复工作。} 题面模型：快慢指针在有环时必然相遇，无环时快指针先到空。暴力版的慢点要落到本题：慢版本会围绕“快慢指针在有环时必然相遇，无环时快指针先到空”借助数组或多次扫描确认目标顺序。样例规律是“规律是快慢指针的相遇说明有环；无环时快指针会先到空”，说明优化点不在节点值，而在少量指针角色能否保持已处理链、未处理链和接回位置都不丢。后面的优化只消掉这类重复工作，不能改变答案集合。

\textbf{优化条件。} 本题的关键观察是：不用哈希表也能常数空间检测环。这句话必须能翻译成可维护状态；如果题目条件改变后观察不再成立，就不能继续套原来的写法。

\textbf{相邻状态。} 规律是快慢指针的相遇说明有环；无环时快指针会先到空。把这个特例继续拆开看：每个指针都要有角色，上一段尾、当前段头、当前节点和后继入口不能混。只要一次改 next 之前没有保存后继，就可能丢掉整段未处理链。

\textbf{状态复用。} 把数组辅助结构换成几个稳定指针角色；重连顺序必须保证已处理链合法、未处理链仍可达。本题落点是：规律是快慢指针的相遇说明有环；无环时快指针会先到空。手算时只改被新输入影响的那一小部分，而不是回头重算完整候选。

\textbf{指针和字段。} 状态字段要能说清：虚拟头、上一段尾、当前段头、当前节点、后继节点；任何改 next 的操作之前都要保存后续入口。手算时按这个协议跟踪：拿 1->2->3 且 3 指回 2 手算，快指针每次走两步，慢指针每次走一步，进入环后二者距离会被不断缩小直到相遇。然后按协议复查：手算时画出 prev、cur、next 和下一段头节点。每次改 next 前先保存后继，这是链表推导的安全线。

\textbf{代码落点。} 所有重连步骤按保存后继、断开或反转、接回前缀、接回后缀的顺序写；最后检查是否形成环或丢节点。关键代码行必须能逐句翻译成“旧状态怎样变成新状态”，否则就只是模板。

\textbf{正确性。} 证明从链结构开始：已处理链连通、未处理链可达、二者连接点明确，节点没有丢失也没有成环。

\textbf{边界实验。} 先用空链、单节点自环、两节点环、尾部无环做反例检查。实验时覆盖空链、单节点自环、两节点环、尾部无环，每步画出前驱、当前、后继和下一段头。链表题不要只看输出值，还要检查节点是否丢失或成环。

\textbf{复杂度变化。} 暴力可能多次扫描链表或拷贝到数组；优化后每个节点通常只被访问、重连常数次，目标是线性时间和常数额外空间。

\textbf{迁移追问。} 如果题目改成“若要求环入口，用相遇后双指针同步走”，先判断原状态是否仍然覆盖新答案集合，再决定是微调边界、改 value 语义，还是换算法。

\noindent\textbf{LeetCode 142 环形链表 II。}

\textbf{题目说明。} LeetCode 142 环形链表 II 的题面入口要先还原成具体任务：相遇点到入口的距离关系来自快慢指针速度差。

\textbf{样例入口。} 拿 1->2->3->4 且 4 指回 2 手算，快慢指针在环内相遇后，把一个指针放回头节点，两个指针同速走，会在入环点 2 相遇。

\textbf{暴力基线。} 慢版本可以先把节点顺序抄到数组里，按数组推导目标顺序。回到链表后，难点不再是值的顺序，而是节点身份和 next 指针不能丢。放到本题就是：先按“相遇点到入口的距离关系来自快慢指针速度差”完整试慢版本；样例里已经能看到“规律是相遇点到入环点的距离和头节点到入环点的距离同余”，所以优化前必须先能说清慢版本枚举了哪些候选。

\textbf{暴力过程。} 拿 1->2->3->4 且 4 指回 2 手算，快慢指针在环内相遇后，把一个指针放回头节点，两个指针同速走，会在入环点 2 相遇。

\textbf{重复工作。} 题面模型：相遇点到入口的距离关系来自快慢指针速度差。暴力版的慢点要落到本题：慢版本会围绕“相遇点到入口的距离关系来自快慢指针速度差”借助数组或多次扫描确认目标顺序。样例规律是“规律是相遇点到入环点的距离和头节点到入环点的距离同余”，说明优化点不在节点值，而在少量指针角色能否保持已处理链、未处理链和接回位置都不丢。后面的优化只消掉这类重复工作，不能改变答案集合。

\textbf{优化条件。} 本题的关键观察是：相遇后一个指针回头，两个指针每次走一步，再次相遇就是入口。这句话必须能翻译成可维护状态；如果题目条件改变后观察不再成立，就不能继续套原来的写法。

\textbf{相邻状态。} 规律是相遇点到入环点的距离和头节点到入环点的距离同余。把这个特例继续拆开看：每个指针都要有角色，上一段尾、当前段头、当前节点和后继入口不能混。只要一次改 next 之前没有保存后继，就可能丢掉整段未处理链。

\textbf{状态复用。} 把数组辅助结构换成几个稳定指针角色；重连顺序必须保证已处理链合法、未处理链仍可达。本题落点是：规律是相遇点到入环点的距离和头节点到入环点的距离同余。手算时只改被新输入影响的那一小部分，而不是回头重算完整候选。

\textbf{指针和字段。} 状态字段要能说清：虚拟头、上一段尾、当前段头、当前节点、后继节点；任何改 next 的操作之前都要保存后续入口。手算时按这个协议跟踪：拿 1->2->3->4 且 4 指回 2 手算，快慢指针在环内相遇后，把一个指针放回头节点，两个指针同速走，会在入环点 2 相遇。然后按协议复查：手算时画出 prev、cur、next 和下一段头节点。每次改 next 前先保存后继，这是链表推导的安全线。

\textbf{代码落点。} 所有重连步骤按保存后继、断开或反转、接回前缀、接回后缀的顺序写；最后检查是否形成环或丢节点。关键代码行必须能逐句翻译成“旧状态怎样变成新状态”，否则就只是模板。

\textbf{正确性。} 证明从链结构开始：已处理链连通、未处理链可达、二者连接点明确，节点没有丢失也没有成环。

\textbf{边界实验。} 先用入口在头节点、长尾短环、无环、单节点环做反例检查。实验时覆盖入口在头节点、长尾短环、无环、单节点环，每步画出前驱、当前、后继和下一段头。链表题不要只看输出值，还要检查节点是否丢失或成环。

\textbf{复杂度变化。} 暴力可能多次扫描链表或拷贝到数组；优化后每个节点通常只被访问、重连常数次，目标是线性时间和常数额外空间。

\textbf{迁移追问。} 如果题目改成“也可用哈希集合，但空间更高”，先判断原状态是否仍然覆盖新答案集合，再决定是微调边界、改 value 语义，还是换算法。

\noindent\textbf{LeetCode 146 LRU 缓存。}

\textbf{题目说明。} LeetCode 146 LRU 缓存 的题面入口要先还原成具体任务：哈希表负责按 key 定位，双向链表负责维护新旧顺序。

\textbf{样例入口。} 拿容量 2 依次 put(1)、put(2)、get(1)、put(3) 手算。get(1) 后 1 变成最新，淘汰时应删 2。数组维护顺序每次移动成本高；链表移动节点到头 O(1)，哈希表按 key 定位节点 O(1)。

\textbf{暴力基线。} 暴力用数组保存缓存项，每次 get 线性查找 key，命中后把元素移动到最新位置；put 时线性找 key 或淘汰最旧项。

\textbf{暴力过程。} 容量 2，put(1)、put(2)、get(1)、put(3) 时，get(1) 会让 1 变成最新，put(3) 应淘汰最旧的 2。

\textbf{重复工作。} 题面模型：哈希表负责按 key 定位，双向链表负责维护新旧顺序。暴力版的慢点要落到本题：浪费在按 key 查找和移动新旧顺序都线性。LRU 需要两个能力：O(1) 按 key 定位节点，O(1) 把节点移动到头部和删除尾部。后面的优化只消掉这类重复工作，不能改变答案集合。

\textbf{优化条件。} 本题的关键观察是：访问和更新都把节点移动到头部，容量满时删除尾部最旧节点。这句话必须能翻译成可维护状态；如果题目条件改变后观察不再成立，就不能继续套原来的写法。

\textbf{相邻状态。} 规律是哈希负责找节点，双向链表负责维护新旧顺序。把这个特例继续拆开看：每个指针都要有角色，上一段尾、当前段头、当前节点和后继入口不能混。只要一次改 next 之前没有保存后继，就可能丢掉整段未处理链。

\textbf{状态复用。} 哈希表 key->链表节点迭代器，双向链表按新到旧保存 key/value。访问或更新节点时移动到头部，容量超限删除尾部。手算时只改被新输入影响的那一小部分，而不是回头重算完整候选。

\textbf{指针和字段。} 状态字段要能说清：map 定位缓存项，list.front 是最新项，list.back 是最旧项，capacity 控制淘汰边界。手算时按这个协议跟踪：get(1) 命中后，把 1 对应节点移到表头；put(3) 超容量时，删除表尾 2，并从哈希表同步 erase。

\textbf{代码落点。} C++ 可用 std::list<pair<int, int>> 保存顺序，unordered\_map<int, list::iterator> 定位。更新已有 key 时先改值再 splice 到头部。关键代码行必须能逐句翻译成“旧状态怎样变成新状态”，否则就只是模板。

\textbf{正确性。} 每次访问都把对应项放到最新位置，所以链表尾部始终是最长时间未访问项；哈希表和链表同步维护同一批 key。

\textbf{边界实验。} 先用容量为一、重复 put、get 不存在、更新已有值做反例检查。实验时覆盖容量为一、重复 put、get 不存在、更新已有值，每步画出前驱、当前、后继和下一段头。链表题不要只看输出值，还要检查节点是否丢失或成环。

\textbf{复杂度变化。} 数组暴力 get/put O(n)；哈希表加双链表 get/put 平均 O(1)，空间 O(capacity)。

\textbf{迁移追问。} 如果题目改成“C++ 可用 \code{std::list} 加迭代器，也可手写双向链表理解所有权”，先判断原状态是否仍然覆盖新答案集合，再决定是微调边界、改 value 语义，还是换算法。

\noindent\textbf{LeetCode 148 排序链表。}

\textbf{题目说明。} LeetCode 148 排序链表 的题面入口要先还原成具体任务：链表不能随机访问，适合归并排序而不是快速排序数组 partition。

\textbf{样例入口。} 拿 4->2->1->3 手算，快慢指针把链表切成 4->2 和 1->3，分别排好后再归并。

\textbf{暴力基线。} 慢版本可以先把节点顺序抄到数组里，按数组推导目标顺序。回到链表后，难点不再是值的顺序，而是节点身份和 next 指针不能丢。放到本题就是：先按“链表不能随机访问，适合归并排序而不是快速排序数组 partition”完整试慢版本；样例里已经能看到“规律是链表排序不能随机访问，归并排序只需要切断中点和有序合并，复杂度稳定在 O(n log n)”，所以优化前必须先能说清慢版本枚举了哪些候选。

\textbf{暴力过程。} 拿 4->2->1->3 手算，快慢指针把链表切成 4->2 和 1->3，分别排好后再归并。

\textbf{重复工作。} 题面模型：链表不能随机访问，适合归并排序而不是快速排序数组 partition。暴力版的慢点要落到本题：慢版本会围绕“链表不能随机访问，适合归并排序而不是快速排序数组 partition”借助数组或多次扫描确认目标顺序。样例规律是“规律是链表排序不能随机访问，归并排序只需要切断中点和有序合并，复杂度稳定在 O(n log n)”，说明优化点不在节点值，而在少量指针角色能否保持已处理链、未处理链和接回位置都不丢。后面的优化只消掉这类重复工作，不能改变答案集合。

\textbf{优化条件。} 本题的关键观察是：自底向上迭代归并能避免递归栈，同时保持 O(n log n)。这句话必须能翻译成可维护状态；如果题目条件改变后观察不再成立，就不能继续套原来的写法。

\textbf{相邻状态。} 规律是链表排序不能随机访问，归并排序只需要切断中点和有序合并，复杂度稳定在 O(n log n)。把这个特例继续拆开看：每个指针都要有角色，上一段尾、当前段头、当前节点和后继入口不能混。只要一次改 next 之前没有保存后继，就可能丢掉整段未处理链。

\textbf{状态复用。} 把数组辅助结构换成几个稳定指针角色；重连顺序必须保证已处理链合法、未处理链仍可达。本题落点是：规律是链表排序不能随机访问，归并排序只需要切断中点和有序合并，复杂度稳定在 O(n log n)。手算时只改被新输入影响的那一小部分，而不是回头重算完整候选。

\textbf{指针和字段。} 状态字段要能说清：虚拟头、上一段尾、当前段头、当前节点、后继节点；任何改 next 的操作之前都要保存后续入口。手算时按这个协议跟踪：拿 4->2->1->3 手算，快慢指针把链表切成 4->2 和 1->3，分别排好后再归并。然后按协议复查：手算时画出 prev、cur、next 和下一段头节点。每次改 next 前先保存后继，这是链表推导的安全线。

\textbf{代码落点。} 所有重连步骤按保存后继、断开或反转、接回前缀、接回后缀的顺序写；最后检查是否形成环或丢节点。关键代码行必须能逐句翻译成“旧状态怎样变成新状态”，否则就只是模板。

\textbf{正确性。} 证明从链结构开始：已处理链连通、未处理链可达、二者连接点明确，节点没有丢失也没有成环。

\textbf{边界实验。} 先用空链、单节点、重复值、奇数长度、切断子链做反例检查。实验时覆盖空链、单节点、重复值、奇数长度、切断子链，每步画出前驱、当前、后继和下一段头。链表题不要只看输出值，还要检查节点是否丢失或成环。

\textbf{复杂度变化。} 暴力可能多次扫描链表或拷贝到数组；优化后每个节点通常只被访问、重连常数次，目标是线性时间和常数额外空间。

\textbf{迁移追问。} 如果题目改成“若把值复制到数组排序，会改变节点身份语义”，先判断原状态是否仍然覆盖新答案集合，再决定是微调边界、改 value 语义，还是换算法。

\noindent\textbf{LeetCode 160 相交链表。}

\textbf{题目说明。} LeetCode 160 相交链表 的题面入口要先还原成具体任务：相交是节点地址相同，不是节点值相同。

\textbf{样例入口。} 拿 A: a1->a2->c1->c2 和 B: b1->c1->c2 手算，交点是节点地址 c1，不是值相等。两个指针走到尾后切换到对方头部，会在总路程相同后于 c1 相遇。

\textbf{暴力基线。} 慢版本可以先把节点顺序抄到数组里，按数组推导目标顺序。回到链表后，难点不再是值的顺序，而是节点身份和 next 指针不能丢。放到本题就是：先按“相交是节点地址相同，不是节点值相同”完整试慢版本；样例里已经能看到“规律是比较节点身份，长度差由切换链表自动抵消”，所以优化前必须先能说清慢版本枚举了哪些候选。

\textbf{暴力过程。} 拿 A: a1->a2->c1->c2 和 B: b1->c1->c2 手算，交点是节点地址 c1，不是值相等。两个指针走到尾后切换到对方头部，会在总路程相同后于 c1 相遇。

\textbf{重复工作。} 题面模型：相交是节点地址相同，不是节点值相同。暴力版的慢点要落到本题：慢版本会围绕“相交是节点地址相同，不是节点值相同”借助数组或多次扫描确认目标顺序。样例规律是“规律是比较节点身份，长度差由切换链表自动抵消”，说明优化点不在节点值，而在少量指针角色能否保持已处理链、未处理链和接回位置都不丢。后面的优化只消掉这类重复工作，不能改变答案集合。

\textbf{优化条件。} 本题的关键观察是：两个指针走完各自链表后切换到另一条链，走过总长度相同后相遇。这句话必须能翻译成可维护状态；如果题目条件改变后观察不再成立，就不能继续套原来的写法。

\textbf{相邻状态。} 规律是比较节点身份，长度差由切换链表自动抵消。把这个特例继续拆开看：每个指针都要有角色，上一段尾、当前段头、当前节点和后继入口不能混。只要一次改 next 之前没有保存后继，就可能丢掉整段未处理链。

\textbf{状态复用。} 把数组辅助结构换成几个稳定指针角色；重连顺序必须保证已处理链合法、未处理链仍可达。本题落点是：规律是比较节点身份，长度差由切换链表自动抵消。手算时只改被新输入影响的那一小部分，而不是回头重算完整候选。

\textbf{指针和字段。} 状态字段要能说清：虚拟头、上一段尾、当前段头、当前节点、后继节点；任何改 next 的操作之前都要保存后续入口。手算时按这个协议跟踪：拿 A: a1->a2->c1->c2 和 B: b1->c1->c2 手算，交点是节点地址 c1，不是值相等。两个指针走到尾后切换到对方头部，会在总路程相同后于 c1 相遇。然后按协议复查：手算时画出 prev、cur、next 和下一段头节点。每次改 next 前先保存后继，这是链表推导的安全线。

\textbf{代码落点。} 所有重连步骤按保存后继、断开或反转、接回前缀、接回后缀的顺序写；最后检查是否形成环或丢节点。关键代码行必须能逐句翻译成“旧状态怎样变成新状态”，否则就只是模板。

\textbf{正确性。} 证明从链结构开始：已处理链连通、未处理链可达、二者连接点明确，节点没有丢失也没有成环。

\textbf{边界实验。} 先用不相交、头节点相交、尾节点相交、长度差很大做反例检查。实验时覆盖不相交、头节点相交、尾节点相交、长度差很大，每步画出前驱、当前、后继和下一段头。链表题不要只看输出值，还要检查节点是否丢失或成环。

\textbf{复杂度变化。} 暴力可能多次扫描链表或拷贝到数组；优化后每个节点通常只被访问、重连常数次，目标是线性时间和常数额外空间。

\textbf{迁移追问。} 如果题目改成“哈希集合更直观但空间更高”，先判断原状态是否仍然覆盖新答案集合，再决定是微调边界、改 value 语义，还是换算法。

\noindent\textbf{LeetCode 206 反转链表。}

\textbf{题目说明。} LeetCode 206 反转链表 的题面入口要先还原成具体任务：每次把当前节点的 next 指向已经反转好的前缀头。

\textbf{样例入口。} 拿 1->2->3 手算，处理 1 时先保存 next=2，再让 1->null；处理 2 时保存 3，再让 2->1。

\textbf{暴力基线。} 慢版本可以先把节点顺序抄到数组里，按数组推导目标顺序。回到链表后，难点不再是值的顺序，而是节点身份和 next 指针不能丢。放到本题就是：先按“每次把当前节点的 next 指向已经反转好的前缀头”完整试慢版本；样例里已经能看到“规律是 prev 保存已反转前缀头，cur 保存待处理节点，每次改 next 前必须保存后继”，所以优化前必须先能说清慢版本枚举了哪些候选。

\textbf{暴力过程。} 拿 1->2->3 手算，处理 1 时先保存 next=2，再让 1->null；处理 2 时保存 3，再让 2->1。

\textbf{重复工作。} 题面模型：每次把当前节点的 next 指向已经反转好的前缀头。暴力版的慢点要落到本题：慢版本会围绕“每次把当前节点的 next 指向已经反转好的前缀头”借助数组或多次扫描确认目标顺序。样例规律是“规律是 prev 保存已反转前缀头，cur 保存待处理节点，每次改 next 前必须保存后继”，说明优化点不在节点值，而在少量指针角色能否保持已处理链、未处理链和接回位置都不丢。后面的优化只消掉这类重复工作，不能改变答案集合。

\textbf{优化条件。} 本题的关键观察是：先保存后继，再改指针，再推进 previous 和 current。这句话必须能翻译成可维护状态；如果题目条件改变后观察不再成立，就不能继续套原来的写法。

\textbf{相邻状态。} 规律是 prev 保存已反转前缀头，cur 保存待处理节点，每次改 next 前必须保存后继。把这个特例继续拆开看：每个指针都要有角色，上一段尾、当前段头、当前节点和后继入口不能混。只要一次改 next 之前没有保存后继，就可能丢掉整段未处理链。

\textbf{状态复用。} 把数组辅助结构换成几个稳定指针角色；重连顺序必须保证已处理链合法、未处理链仍可达。本题落点是：规律是 prev 保存已反转前缀头，cur 保存待处理节点，每次改 next 前必须保存后继。手算时只改被新输入影响的那一小部分，而不是回头重算完整候选。

\textbf{指针和字段。} 状态字段要能说清：虚拟头、上一段尾、当前段头、当前节点、后继节点；任何改 next 的操作之前都要保存后续入口。手算时按这个协议跟踪：拿 1->2->3 手算，处理 1 时先保存 next=2，再让 1->null；处理 2 时保存 3，再让 2->1。然后按协议复查：手算时画出 prev、cur、next 和下一段头节点。每次改 next 前先保存后继，这是链表推导的安全线。

\textbf{代码落点。} 所有重连步骤按保存后继、断开或反转、接回前缀、接回后缀的顺序写；最后检查是否形成环或丢节点。关键代码行必须能逐句翻译成“旧状态怎样变成新状态”，否则就只是模板。

\textbf{正确性。} 证明从链结构开始：已处理链连通、未处理链可达、二者连接点明确，节点没有丢失也没有成环。

\textbf{边界实验。} 先用空链、单节点、两个节点、不要丢失剩余链做反例检查。实验时覆盖空链、单节点、两个节点、不要丢失剩余链，每步画出前驱、当前、后继和下一段头。链表题不要只看输出值，还要检查节点是否丢失或成环。

\textbf{复杂度变化。} 暴力可能多次扫描链表或拷贝到数组；优化后每个节点通常只被访问、重连常数次，目标是线性时间和常数额外空间。

\textbf{迁移追问。} 如果题目改成“K 组反转是在这段局部反转外增加分组边界”，先判断原状态是否仍然覆盖新答案集合，再决定是微调边界、改 value 语义，还是换算法。

\noindent\textbf{LeetCode 234 回文链表。}

\textbf{题目说明。} LeetCode 234 回文链表 的题面入口要先还原成具体任务：链表回文要求前半段和后半段逆序后逐一相等。

\textbf{样例入口。} 拿 1->2->2->1 手算，快慢指针找到后半段起点，再反转后半段得到 1->2，与前半段逐个比较。

\textbf{暴力基线。} 慢版本可以先把节点顺序抄到数组里，按数组推导目标顺序。回到链表后，难点不再是值的顺序，而是节点身份和 next 指针不能丢。放到本题就是：先按“链表回文要求前半段和后半段逆序后逐一相等”完整试慢版本；样例里已经能看到“规律是链表不能随机访问，先找中点、反转后半、比较，再按需要恢复结构”，所以优化前必须先能说清慢版本枚举了哪些候选。

\textbf{暴力过程。} 拿 1->2->2->1 手算，快慢指针找到后半段起点，再反转后半段得到 1->2，与前半段逐个比较。

\textbf{重复工作。} 题面模型：链表回文要求前半段和后半段逆序后逐一相等。暴力版的慢点要落到本题：慢版本会围绕“链表回文要求前半段和后半段逆序后逐一相等”借助数组或多次扫描确认目标顺序。样例规律是“规律是链表不能随机访问，先找中点、反转后半、比较，再按需要恢复结构”，说明优化点不在节点值，而在少量指针角色能否保持已处理链、未处理链和接回位置都不丢。后面的优化只消掉这类重复工作，不能改变答案集合。

\textbf{优化条件。} 本题的关键观察是：快慢指针找中点，反转后半段，再和前半段同步比较。这句话必须能翻译成可维护状态；如果题目条件改变后观察不再成立，就不能继续套原来的写法。

\textbf{相邻状态。} 规律是链表不能随机访问，先找中点、反转后半、比较，再按需要恢复结构。把这个特例继续拆开看：每个指针都要有角色，上一段尾、当前段头、当前节点和后继入口不能混。只要一次改 next 之前没有保存后继，就可能丢掉整段未处理链。

\textbf{状态复用。} 把数组辅助结构换成几个稳定指针角色；重连顺序必须保证已处理链合法、未处理链仍可达。本题落点是：规律是链表不能随机访问，先找中点、反转后半、比较，再按需要恢复结构。手算时只改被新输入影响的那一小部分，而不是回头重算完整候选。

\textbf{指针和字段。} 状态字段要能说清：虚拟头、上一段尾、当前段头、当前节点、后继节点；任何改 next 的操作之前都要保存后续入口。手算时按这个协议跟踪：拿 1->2->2->1 手算，快慢指针找到后半段起点，再反转后半段得到 1->2，与前半段逐个比较。然后按协议复查：手算时画出 prev、cur、next 和下一段头节点。每次改 next 前先保存后继，这是链表推导的安全线。

\textbf{代码落点。} 所有重连步骤按保存后继、断开或反转、接回前缀、接回后缀的顺序写；最后检查是否形成环或丢节点。关键代码行必须能逐句翻译成“旧状态怎样变成新状态”，否则就只是模板。

\textbf{正确性。} 证明从链结构开始：已处理链连通、未处理链可达、二者连接点明确，节点没有丢失也没有成环。

\textbf{边界实验。} 先用空链、单节点、奇数长度、偶数长度、是否恢复链表做反例检查。实验时覆盖空链、单节点、奇数长度、偶数长度、是否恢复链表，每步画出前驱、当前、后继和下一段头。链表题不要只看输出值，还要检查节点是否丢失或成环。

\textbf{复杂度变化。} 暴力可能多次扫描链表或拷贝到数组；优化后每个节点通常只被访问、重连常数次，目标是线性时间和常数额外空间。

\textbf{迁移追问。} 如果题目改成“如果不允许修改链表，可用栈保存前半段但空间变为线性”，先判断原状态是否仍然覆盖新答案集合，再决定是微调边界、改 value 语义，还是换算法。

\noindent\textbf{LeetCode 430 扁平化多级双向链表。}

\textbf{题目说明。} LeetCode 430 扁平化多级双向链表 的题面入口要先还原成具体任务：多级链表按深度优先顺序展开成一条双向链表。

\textbf{样例入口。} 拿 1-2-3 且 2 有子链 4-5 手算，扁平化顺序应为 1-2-4-5-3；若不保存 2 原来的 next=3，子链接完后会丢失 3。

\textbf{暴力基线。} 慢版本可以先把节点顺序抄到数组里，按数组推导目标顺序。回到链表后，难点不再是值的顺序，而是节点身份和 next 指针不能丢。放到本题就是：先按“多级链表按深度优先顺序展开成一条双向链表”完整试慢版本；样例里已经能看到“规律是遇到 child 时先保存 next，再把 child 接入，并把 child 尾部接回原 next”，所以优化前必须先能说清慢版本枚举了哪些候选。

\textbf{暴力过程。} 拿 1-2-3 且 2 有子链 4-5 手算，扁平化顺序应为 1-2-4-5-3；若不保存 2 原来的 next=3，子链接完后会丢失 3。

\textbf{重复工作。} 题面模型：多级链表按深度优先顺序展开成一条双向链表。暴力版的慢点要落到本题：慢版本会围绕“多级链表按深度优先顺序展开成一条双向链表”借助数组或多次扫描确认目标顺序。样例规律是“规律是遇到 child 时先保存 next，再把 child 接入，并把 child 尾部接回原 next”，说明优化点不在节点值，而在少量指针角色能否保持已处理链、未处理链和接回位置都不丢。后面的优化只消掉这类重复工作，不能改变答案集合。

\textbf{优化条件。} 本题的关键观察是：显式栈保存后续 next，遇到 child 先接 child，再在子链结束后恢复原 next。这句话必须能翻译成可维护状态；如果题目条件改变后观察不再成立，就不能继续套原来的写法。

\textbf{相邻状态。} 规律是遇到 child 时先保存 next，再把 child 接入，并把 child 尾部接回原 next。把这个特例继续拆开看：每个指针都要有角色，上一段尾、当前段头、当前节点和后继入口不能混。只要一次改 next 之前没有保存后继，就可能丢掉整段未处理链。

\textbf{状态复用。} 把数组辅助结构换成几个稳定指针角色；重连顺序必须保证已处理链合法、未处理链仍可达。本题落点是：规律是遇到 child 时先保存 next，再把 child 接入，并把 child 尾部接回原 next。手算时只改被新输入影响的那一小部分，而不是回头重算完整候选。

\textbf{指针和字段。} 状态字段要能说清：虚拟头、上一段尾、当前段头、当前节点、后继节点；任何改 next 的操作之前都要保存后续入口。手算时按这个协议跟踪：拿 1-2-3 且 2 有子链 4-5 手算，扁平化顺序应为 1-2-4-5-3；若不保存 2 原来的 next=3，子链接完后会丢失 3。然后按协议复查：手算时画出 prev、cur、next 和下一段头节点。每次改 next 前先保存后继，这是链表推导的安全线。

\textbf{代码落点。} 所有重连步骤按保存后继、断开或反转、接回前缀、接回后缀的顺序写；最后检查是否形成环或丢节点。关键代码行必须能逐句翻译成“旧状态怎样变成新状态”，否则就只是模板。

\textbf{正确性。} 证明从链结构开始：已处理链连通、未处理链可达、二者连接点明确，节点没有丢失也没有成环。

\textbf{边界实验。} 先用child 为空、next 和 child 同时存在、多层嵌套、prev 指针恢复做反例检查。实验时覆盖child 为空、next 和 child 同时存在、多层嵌套、prev 指针恢复，每步画出前驱、当前、后继和下一段头。链表题不要只看输出值，还要检查节点是否丢失或成环。

\textbf{复杂度变化。} 暴力可能多次扫描链表或拷贝到数组；优化后每个节点通常只被访问、重连常数次，目标是线性时间和常数额外空间。

\textbf{迁移追问。} 如果题目改成“它和二叉树前序展开类似，核心是保存未处理分支”，先判断原状态是否仍然覆盖新答案集合，再决定是微调边界、改 value 语义，还是换算法。

\noindent\textbf{LeetCode 460 LFU 缓存。}

\textbf{题目说明。} LeetCode 460 LFU 缓存 的题面入口要先还原成具体任务：缓存淘汰同时按访问频次和同频最近性排序，需要维护两个维度的索引。

\textbf{样例入口。} 拿容量 2，put(1)、put(2)、get(1)、put(3) 手算。此时 1 的频次高于 2，应淘汰 2；若两个 key 频次相同，还要淘汰同频里最旧的。

\textbf{暴力基线。} 慢版本可以先把节点顺序抄到数组里，按数组推导目标顺序。回到链表后，难点不再是值的顺序，而是节点身份和 next 指针不能丢。放到本题就是：先按“缓存淘汰同时按访问频次和同频最近性排序，需要维护两个维度的索引”完整试慢版本；样例里已经能看到“规律是 key 哈希定位节点，频次哈希定位同频链表，minFreq 指向当前最低频次”，所以优化前必须先能说清慢版本枚举了哪些候选。

\textbf{暴力过程。} 拿容量 2，put(1)、put(2)、get(1)、put(3) 手算。此时 1 的频次高于 2，应淘汰 2；若两个 key 频次相同，还要淘汰同频里最旧的。

\textbf{重复工作。} 题面模型：缓存淘汰同时按访问频次和同频最近性排序，需要维护两个维度的索引。暴力版的慢点要落到本题：慢版本会围绕“缓存淘汰同时按访问频次和同频最近性排序，需要维护两个维度的索引”借助数组或多次扫描确认目标顺序。样例规律是“规律是 key 哈希定位节点，频次哈希定位同频链表，minFreq 指向当前最低频次”，说明优化点不在节点值，而在少量指针角色能否保持已处理链、未处理链和接回位置都不丢。后面的优化只消掉这类重复工作，不能改变答案集合。

\textbf{优化条件。} 本题的关键观察是：哈希表按 key 定位节点，频次到双向链表保存同频节点顺序，并维护当前最小频次。这句话必须能翻译成可维护状态；如果题目条件改变后观察不再成立，就不能继续套原来的写法。

\textbf{相邻状态。} 规律是 key 哈希定位节点，频次哈希定位同频链表，minFreq 指向当前最低频次。把这个特例继续拆开看：每个指针都要有角色，上一段尾、当前段头、当前节点和后继入口不能混。只要一次改 next 之前没有保存后继，就可能丢掉整段未处理链。

\textbf{状态复用。} 把数组辅助结构换成几个稳定指针角色；重连顺序必须保证已处理链合法、未处理链仍可达。本题落点是：规律是 key 哈希定位节点，频次哈希定位同频链表，minFreq 指向当前最低频次。手算时只改被新输入影响的那一小部分，而不是回头重算完整候选。

\textbf{指针和字段。} 状态字段要能说清：虚拟头、上一段尾、当前段头、当前节点、后继节点；任何改 next 的操作之前都要保存后续入口。手算时按这个协议跟踪：拿容量 2，put(1)、put(2)、get(1)、put(3) 手算。此时 1 的频次高于 2，应淘汰 2；若两个 key 频次相同，还要淘汰同频里最旧的。然后按协议复查：手算时画出 prev、cur、next 和下一段头节点。每次改 next 前先保存后继，这是链表推导的安全线。

\textbf{代码落点。} 所有重连步骤按保存后继、断开或反转、接回前缀、接回后缀的顺序写；最后检查是否形成环或丢节点。关键代码行必须能逐句翻译成“旧状态怎样变成新状态”，否则就只是模板。

\textbf{正确性。} 证明从链结构开始：已处理链连通、未处理链可达、二者连接点明确，节点没有丢失也没有成环。

\textbf{边界实验。} 先用容量为零、更新已有 key、频次链表变空时最小频次更新、同频淘汰最旧节点做反例检查。实验时覆盖容量为零、更新已有 key、频次链表变空时最小频次更新、同频淘汰最旧节点，每步画出前驱、当前、后继和下一段头。链表题不要只看输出值，还要检查节点是否丢失或成环。

\textbf{复杂度变化。} 暴力可能多次扫描链表或拷贝到数组；优化后每个节点通常只被访问、重连常数次，目标是线性时间和常数额外空间。

\textbf{迁移追问。} 如果题目改成“LRU 只有时间顺序一个维度，LFU 多了频次维度，结构维护明显更复杂”，先判断原状态是否仍然覆盖新答案集合，再决定是微调边界、改 value 语义，还是换算法。

\subsection{实验复盘}

追问通常会让你画指针。请准备说明上一段尾、当前段头、当前节点、后继节点分别是谁，以及哪一步保存 next 防止断链。

复盘本节时先从 LeetCode 2 两数相加、LeetCode 19 删除链表的倒数第 N 个结点、LeetCode 21 合并两个有序链表 开始：每道题都先手写一个能覆盖全部候选的暴力版，再指出优化结构具体保存了哪一段历史信息，最后用相邻案例比较为什么不能互换解法。
